\section{Related Work}
\label{sec:background}

Machine unlearning in diffusion models aims to remove a target concept $c$ (e.g., \textit{horse}) from a pretrained model $\mathcal{D}_\theta$
% \varun{use a different variable for params, as the params $\theta$ are used for the U-Net?} 
% \yian{$\theta$ is shared with $\epsilon_\theta$ on purpose as the U-Net is the only trainable component, so $\mathcal{D}_\theta$ and $\epsilon_\theta$ must share parameters. }
while preserving overall generative capabilities. The goal is an updated model $\mathcal{D}_{\hat{\theta}}$ that suppresses $c$ on target-associated prompts while maintaining quality on unrelated prompts, following prior work on concept erasure in diffusion models. Background on latent diffusion models is in Appendix~\ref{app:background}.

\paragraph{Closed-form editing.} UCE~\citep{gandikota2024unified} and RECE~\citep{gong2024reliable} modify cross-attention projections through analytic updates, without gradient training. RECE additionally searches for embeddings that regenerate the target and removes them in closed form.

\textbf{Fine-tuning methods.}
A large body of work erase concepts by directly modifying the weights of a pretrained diffusion model. Erased Stable Diffusion (ESD)~\cite{gandikota2023erasing} uses negative guidance to align target-token prompts with a neutral anchor. Subsequent work improves the precision by editing cross-attention representations~\citep{wang2025ace} or updating salient concept-related parameters~\citep{fan2023salun,wu2024scissorhands}. CPE~\citep{lee2025concept} trains residual attention gates with an anchoring loss, and TRUST~\citep{kori2026selective} selectively fine-tunes target-associated neurons under Hessian-based regularization. More recent robust methods such as AdvUnlearn~\cite{zhang2024defensive} and STEREO~\cite{srivatsan2025stereo} further defend against prompt-based regeneration attacks.
However, these methods still largely define erasure through target tokens, prompts, or localized parameter subsets. They therefore suppress direct target prompts well, but struggle with \emph{concept leakage} through correlated context; when made more aggressive, their edits can propagate through shared representations and cause \emph{collateral damage} to semantic neighbors.

\textbf{Inference-time methods.}
A complementary line of work avoids modifying $\mathcal{D}_{\theta}$ and instead intervenes during generation by manipulating text embeddings~\citep{li2024get}, projecting prompt tokens away from a target subspace (SAFREE, \citealp{yoon2025safree}), adding safety guidance to the denoising process (SLD, \citealp{schramowski2023safe}), inserting lightweight concept-blocking adapters~\citep{lyu2024one}, or ablating sparse-autoencoder features~\citep{cywinski2025saeuron}. These methods are computationally efficient and often preserve broad model behavior. Intervening at inference time does not by itself make a method conservative: depending on how broadly the intervention displaces activations, these methods span the full range from strong erasure with substantial neighbor damage to strong preservation with residual leakage (Section~\ref{subsec:tradeoff_real}).

\paragraph{Preservation-oriented methods.} Several methods explicitly target the collateral damage we study. CPE anchors non-target concepts during erasure, TRUST restricts updates to target-associated neurons, and SAFREE adapts its filtering to preserve prompt fidelity. Adversarial preservation~\citep{bui2024erasing} identifies and protects the concepts most sensitive to erasure, Receler~\citep{huang2024receler} and MACE~\citep{lu2024mace} localize edits through lightweight erasers and LoRA fusion, and meta-unlearning~\citep{gao2025meta} prevents relearning of erased concepts. We evaluate CPE, TRUST, and SAFREE; they reach more favorable operating points but do not escape the trade-off (Section~\ref{subsec:tradeoff_real}).

\paragraph{Positioning.} That aggressive erasure degrades utility is well
documented, and the preservation-oriented methods above aim to reduce it;
localization work such as \citet{zarei2026localizing} identifies where concepts
are represented. Neither models target-neighbor overlap or derives limits on
unlearning. Our claim is narrower and structural: robust erasure necessarily
degrades related concepts \emph{because} their representations overlap, a limit
we quantify through $\kappa$, prove unavoidable (Theorem~\ref{thm:tradeoff},
Corollary~\ref{cor:impossibility}), and turn into a frontier-based evaluation
protocol. Across all three families, methods implicitly treat concepts as
localized and independently suppressible; this assumption fails whenever
activation regions overlap.

%\ali{we should somehow highlight this last paragraph; readers usually skip the related work sections.}
% \noindent\textbf{Takeaway.}
% Existing unlearning algorithms operate under the assumption that concepts can be selectively removed from the model’s representation space. In the following section, we provide empirical evidence that challenges this assumption by demonstrating that concepts remain recoverable through correlated contextual prompts.
% \noindent \textbf{Data Unlearning}
% \cite{alberti2025data}

\begin{comment}
    \subsection{Unlearning Methods}
% \noindent \textbf{The \textsc{UnlearnCanvas} Benchmark}
% Our experimental framework is built upon the \textsc{UnlearnCanvas} benchmark, a comprehensive dataset recently proposed by~\citet{zhang2024unlearncanvas} to standardize the evaluation of diffusion model unlearning. It was created to address the lack of rigorous, automated evaluation frameworks, which previously relied on ad-hoc, qualitative assessments. The benchmark consists of a high-resolution, stylized image dataset featuring 65 distinct artistic styles and 20 distinct object classes. Its primary advantage lies in its high stylistic consistency, which enables the training of a highly accurate style classifier, allowing for the reliable, quantitative measurement of unlearning using standardized metrics like Unlearning Accuracy (UA), In-domain Retain Accuracy (IRA), and Cross-domain Retain Accuracy (CRA). While \textsc{UnlearnCanvas} provides the essential visual data for our investigation, it does not include a formal, art style hierarchy linking its 60 styles. The styles are presented as a flat list. A foundational part of our methodology, therefore, is to construct this missing knowledge graph, which serves as the basis for our directional unlearning experiments.

\noindent \textbf{Methods for Diffusion Model Unlearning}
The primary objective of Machine Unlearning (MU) in diffusion models is to erase a target concept while minimizing collateral damage to the model's remaining generative capabilities. Current approaches can be broadly classified into two main categories:

\textbf{1. Fine-Tuning Methods:} This line of work involves permanently altering the model's weights. Many of these methods, such as Erased Stable Diffusion (ESD)~\cite{gandikota2023erasing}, utilize a form of negative guidance, retraining the model to align the output of the "forget" prompt with that of a neutral anchor prompt. Variations of this approach aim for greater precision. 
For example,~\citet{zhang2024forget} applies a re-steering loss to attention layers, while~\citet{fan2023salun} and~\citet{wu2024scissorhands} identify and modify only the most sensitive weights. 
\cite{zhang2024defensive} propose \textit{AdvUnlearn}, which integrates adversarial training into diffusion model unlearning to improve robustness against prompt-based attacks that attempt to regenerate erased concepts.
Attention-based Concept Erasure (ACE)~\cite{wang2025ace} removes target concepts by modifying cross-attention representations associated with the concept tokens, enabling localized editing of concept-related activations within the diffusion model.
STEREO~\cite{srivatsan2025stereo} introduces a robust concept unlearning framework that enforces separation between target and neighboring concepts, improving resistance to prompt-based regeneration of erased concepts.
%More recently,~\citet{li2025set} proposes a trajectory-aware finetuning framework that explicitly preserves early-stage denoising dynamics while reversing classifier-free guidance directions during mid-to-late stages, enabling precise concept erasure without relying on heuristic anchor selection.

Other notable fine-tuning methods include Unified Concept Editing (UCE)~\cite{gandikota2024unified}, Concept Ablation~\cite{vinker2023concept}, Selective Amnesia (SA)~\cite{heng2023selective}, and EraseDiff (EDiff)~\cite{wu2024erasediff}. 

While effective, these methods are vulnerable to "jailbreak" attacks that exploit related concepts~\cite{zhang2024generate} and operationally treat concepts as isolated entities. Orthogonal to concept-level unlearning,~\citet{alberti2025data} formalize datapoint-level unlearning in diffusion models and introduce Subtracted Importance Sampled Scores (SISS), an efficient fine-tuning objective with theoretical guarantees that aims to match the retrain-without-data gold standard at far lower cost~\cite{alberti2025data}

\textbf{2. Inference-Time Methods:} A second class of methods avoids the costly fine-tuning step by intervening directly during the generation process, leaving the model's original weights untouched. These techniques are often more efficient and can be dynamically toggled. SEOT~\cite{li2024get}suppresses unwanted content by manipulating text embeddings, and SPM~\cite{lyu2024one} uses lightweight one-dimensional adapters to block the propagation of concept-specific information. A highly relevant new approach in this category is SAeUron~\cite{cywinski2025saeuron}. This method leverages Sparse Autoencoders (SAEs) trained on the model's internal activations to learn a dictionary of sparse, interpretable features.And unlearning is achieved at inference time by identifying the features that correspond to an unwanted concept and ablating them, preventing their contribution to the final image. While these methods also treat concepts in isolation, their interventional nature inspires our work. We leverage similar interventions not just for deletion, but as a method to discover the latent structure of conceptual relationships. 
\end{comment}